\documentclass{article}
\usepackage{graphicx} % Required for inserting images
\usepackage{amsmath}
\usepackage{amsfonts}


\begin{document}

\section*{Ejercicio IV}

Realice lo que se pide, planteando una integral doble y resolviéndola. \\

a) Pruebe que el área del conjunto \( A = \{(x,y) \in \mathbb{R}^2 | x^2 + y^2 \leq 1\} \) es \( \pi \).  \\

b) Pruebe que el área del conjunto \( B = [0, \sqrt{2}] \times [0, \sqrt{2}] \) es 2.

Inciso (a): Área del conjunto \( A \) \\
El conjunto \( A \) es el disco unitario centrado en el origen, es decir,  
\[
A = \{(x,y) \in \mathbb{R}^2 \mid x^2 + y^2 \leq 1\}.
\]  
Para calcular su área, utilizamos el cambio de coordenadas polares demostrado en geometría analítica y empleado en cálculo 2:  
\[
x = r\cos\theta, \quad y = r\sin\theta, \quad r^2 = x^2+y^2
\]  
El elemento de área en coordenadas polares es \( dA = r\, dr\, d\theta \). Como el disco abarca todo el ángulo \( \theta \) desde \( 0 \) hasta \( 2\pi \) y el radio \( r \) varía desde \( 0 \) hasta \( 1 \), podemos aplicar teorema de Fubini (nótese que será análogo para el ejercicio 7 y 8) la integral doble que da el área es:

\[
\iint_A dA = \int_0^{2\pi} \int_0^1 r\, dr\, d\theta.
\]

Resolviendo la integral interna respecto a \( r \):

\[
\int_0^1 r\, dr = \frac{1}{2} r^2 \Big|_0^1 = \frac{1}{2}.
\]

Ahora integramos respecto a \( \theta \):

\[
\int_0^{2\pi} \frac{1}{2} d\theta = \frac{1}{2} \theta \Big|_0^{2\pi} = \frac{1}{2} (2\pi - 0) = \pi.
\]

Por lo tanto, el área del conjunto \( A \) es \( \pi \), como se quería demostrar.\\



sobre la región dada. La población total se obtiene con la integral doble:

\[
P = \int_0^{10} \int_0^{10} 100(y+1) \,dx\, dy.
\]

Como \( f(x,y) \) no depende de \( x \), la integral en \( x \) es sencilla:

   \[
   \int_0^{10} 100(y+1) \,dx = 100(y+1) \int_0^{10} dx = 100(y+1)(10) = 1000(y+1).
   \]

Ahora, integramos en \( y \):

   \[
   P = \int_0^{10} 1000(y+1) \, dy.
   \]

   Resolviendo:

   \[
   P = 1000 \int_0^{10} (y+1) \, dy.
   \]

   La integral es:

   \[
   \int (y+1) \, dy = \frac{y^2}{2} + y.
   \]

   Evaluamos en los límites \( y=0 \) a \( y=10 \):

   \[
   \left[ \frac{y^2}{2} + y \right]_0^{10} = \left( \frac{10^2}{2} + 10 \right) - \left( \frac{0^2}{2} + 0 \right).
   \]

   \[
   = \left( \frac{100}{2} + 10 \right) = (50+10) = 60.
   \]

Ahora solo multiplicamos por 1000:

   \[
   P = 1000 \times 60 = 60000.
   \]

$\therefore$ La población total en la región es 60,000 personas.

\end{document}
